\section{Interactive Hypothesis Correction}
\label{sec:interactive-hypothesis-correction}

SHI exposes two intermediate decisions in multi-event retrieval as explicit
interaction objects. Before retrieval, the \emph{Query Hypothesis} represents
how the input narrative has been interpreted as an ordered event sequence.
After retrieval, \emph{EventTrail} exposes how those events have been aligned
to occurrences within a selected video. Both representations can be inspected
before the user commits a correction.

\begin{figure*}[t]
    \centering
    % \includegraphics[width=\textwidth]{002_hypothesis_interaction.png}
    \fbox{\parbox[c][3.0cm][c]{0.96\linewidth}{\centering
    Hypothesis Interaction figure placeholder}}
    \caption{The two hypothesis-level interactions in SHI.
    \textbf{(a)} Query Hypothesis exposes the inferred event decomposition
    before retrieval and supports direct structural correction.
    \textbf{(b)} EventTrail exposes the decoded temporal hypothesis of a
    retrieved video and previews complete-path alternatives before a correction
    is committed.}
    \label{fig:hypothesis-interaction}
\end{figure*}

\subsection{Query Hypothesis}
\label{sec:query-hypothesis}

Interactive video retrieval systems commonly support query reformulation,
including suggestions derived from intermediate search results and
conversational refinement with relevance feedback
~\cite{lokoc2023queryreformulation,khan2024exquisitor,sharma2025exquisitor}.
Structured interfaces have also been used to express complex and temporal
queries~\cite{wu2021sql,heller2022temporalqueries}. SHI differs in the object
being edited: the system-generated interpretation of the original narrative is
itself exposed as a persistent \emph{Query Hypothesis} before retrieval.

For a multi-event request, SHI constructs an ordered sequence
\[
H_q=(E_1,E_2,\ldots,E_N),
\]
where each automatically generated event is grounded to a verbatim span of the
original query. The original request and its inferred event sequence are shown
together, allowing the user to inspect whether the decomposition preserves the
intended semantic and temporal structure.

Figure~\ref{fig:hypothesis-interaction}(a) shows the Query Hypothesis editor.
Users can edit an event, split it, merge adjacent events, change event order,
or add a new event. Image examples can additionally be associated with
individual events. In contrast to reformulating a single text query and
observing the resulting ranking, these operations directly modify the event
structure that will be used by the downstream retrieval and temporal-alignment
stages.

Edits are previewed before they affect the active hypothesis. The interface
shows the current and proposed event sequences side by side, after which the
user may cancel or apply the modification. Applying a change creates a new
query revision and marks results from the previous revision as outdated; a new
search is started explicitly by the user. Previous revisions can also be
restored through undo. This separates inspection and correction of the query
interpretation from search execution.

\subsection{EventTrail}
\label{sec:event-trail}

Prior systems support temporal query formulation and retrieval over sequences
or temporally related search conditions
~\cite{heller2021explainable,heller2022temporalqueries,khan2026temporalqueries}.
EventTrail instead exposes the \emph{decoded temporal hypothesis} of an
individual retrieved video: the current assignment of each query event to a
specific occurrence in that video.

Figure~\ref{fig:hypothesis-interaction}(b) presents this hypothesis as an
ordered event path, with each event associated with its selected frame and
timestamp. The user can focus on any event and inspect temporally distinct
alternative occurrences under the current constraints.

A central distinction of EventTrail is that an alternative is not presented as
an isolated frame. For each candidate occurrence of the focused event, SHI
decodes a complete chronologically valid path for the entire event sequence.
The user can therefore preview the consequence of selecting an alternative
before changing the active hypothesis. If other event assignments must move to
preserve chronological consistency, those changes are shown together with the
previewed path.

Previewing is non-mutating. The user can \emph{Keep} the current occurrence,
\emph{Use} a previewed alternative, or \emph{Reject} an incorrect temporal
mode. Keeping or using an occurrence introduces an anchor for that event,
whereas rejection excludes the corresponding temporal region. The complete
path is then re-decoded under the updated constraints using the multimodal
evidence retained from the original search.

Because temporal alignment is solved jointly for the complete event sequence,
a local correction may alter assignments of other events. EventTrail makes
these induced changes visible rather than treating each event match
independently. The interaction can be repeated across events, and previous
corrections can be undone when necessary.